\section{Introduction}
Real-time embedded systems demand precise control over timing, concurrency and
resource management.  
Such systems are often used in critical applications where failure to meet
timing constraints can lead to severe consequences. Hence, developing a
real-time system requires careful consideration of the underlying runtime
environment, necessitating a clear understanding of how the system will behave
in order to provide sufficient guarantees on meeting timing requirements. \\
The choice of the runtime environment is a crucial aspect in the design of a
real-time system, as it can significantly impact the system's ability to meet
its desired behaviour through reliable mechanisms it offers. One of the most
popular languages that meets the needs of real-time systems is \textit{Ada},
which has been widely used in the development of safety-critical systems due to
its deterministic constructs that allow high analysability of the code. While
\textit{Ada} remains a robust choice, recent advances in the \textit{Rust}
programming language have introduced new opportunities for real-time systems
development. \\ 
In this paper, a real-time system originally implemented using the \textit{Ada
Ravenscar Profile} is re-implemented in \textit{Rust} using the\textit{Real-Time
Intrupput-driven Concurrency (RTIC)} framework \cite{RTIC}, which is responsible for
managing the scheduling of tasks through a \textit{Fixed-Priority Scheduling}
and the resource sharing in the system under \textit{Stack Resource Policy
(SRP)}. \\ 
For the analysis of the system, the \textit{MAST} tool is used to verify the
timing properties and to analyze the system under different workloads, with the 
system running on a \textit{STM32F4} microcontroller. \\ 
The goal of this work is to demonstrate the feasibility of using \textit{Rust}
and \textit{RTIC} for real-time systems, while also providing some
considerations and comparison with the original \textit{Ada} implementation. \\ 
The rest of the paper is organized as follows. Section 2 presents a formal
description of the system model. Section 3 introduces the \textit{Rust Runtime},
in particular the \textit{RTIC} framework and its key concepts. Sections 4 and 5
discuss the measurement of execution times and the detection of deadline misses,
respectively. Section 6 provides an overview of the \textit{MAST} tool and its
analysis capabilities. In Section 7, we present the \textit{Fixed-Priority
Scheduling (FPS)} analysis of the system. Finally, Section 8 concludes the paper
reporting the main findings.  